\documentclass[11pt]{article}
\usepackage{fontspec,amsmath,amssymb,amsthm,hyperref}
\usepackage[margin=1in]{geometry}
\setlength{\emergencystretch}{3em}
\title{The genuine finite-sum control-variate estimator}
\author{LangevinLean mathematical exposition}
\date{}
\begin{document}\maketitle
\noindent Source module: \texttt{AverageLMC/SGLDEstimator.lean}. Compile this proof document with XeLaTeX or LuaLaTeX. The accompanying HTML pairs the mathematical argument with the exact Lean statements and records the scope of the formalization.


This module constructs the finite-sum potential, the actual uniform component-index probability, and the fixed-point control-variate estimator in the SGLD update. It derives unbiasedness, the exact centered variance identity, the sharp RMS component bound, and all product-space second moments. The anchor is justified by existence and uniqueness of a true minimizer, including its zero gradient.


\section{Construct the potential and the uniform index probability}


The total potential is the finite sum of the component functions. Differentiation commutes with this finite sum, and the gradient identification follows by the inner-product duality map. The uniform index law is counting measure multiplied by 1/n. For n positive its mass is one. It is a finite measure even at n=0, where the index space is empty. Its Bochner integral is exactly the finite vector sum divided by n; this formula is used for all following index moments.



\noindent Lean declarations: \texttt{componentPotential}, \texttt{uniformIndexLaw}, \texttt{uniformIndexLaw\_probability}, \texttt{uniformIndexLaw\_isFiniteMeasure}, \texttt{integral\_uniformIndexLaw}, \texttt{gradient\_componentPotential}.

\section{Use the actual fixed-point estimator and prove unbiasedness}


For an anchor xstar, define the estimator n times (gradient f\_i(x)-gradient f\_i(xstar)). It vanishes at the anchor for every index. Averaging over the uniform index cancels n and sums the component gradients. If the true total gradient at the anchor is zero, this average is exactly gradient f(x). Thus the paper estimator is unbiased for the gradient of the same potential defining its Gibbs target.



\noindent Lean declarations: \texttt{controlVariateEstimator}, \texttt{controlVariateEstimator\_at\_anchor}, \texttt{controlVariateEstimator\_unbiased}.

\section{Control the full gradient by exactly n M\_rms}


Squaring the defining square root gives M\_rms squared equal to the average squared exact component-gradient seminorms. The triangle inequality bounds the gradient Lipschitz constant of the sum by the sum of component constants. Finite Cauchy-Schwarz bounds the sum of exact nonnegative seminorms by n M\_rms. Replacing each supplied component bound by its proved exact seminorm gives a Lipschitz gradient bound for the full potential with exactly n M\_rms. Comparing a strong-convexity lower model with the coordinate upper model then proves m at most n M\_rms. These comparisons justify the dimensionless step-size conditions in the SGLD theorem.



\noindent Lean declarations: \texttt{rmsComponentSmoothness\_sq}, \texttt{componentPotential\_gradient\_lipschitz}, \texttt{sum\_globalSmoothness\_le\_n\_mul\_rms}, \texttt{componentPotential\_gradient\_lipschitz\_rms}, \texttt{strongConvex\_le\_n\_mul\_rms}.

\section{Compute and bound the estimator second moment}


The uniform average of the squared estimator norm equals n times the sum of squared component-gradient differences. Each difference is bounded by the component exact seminorm times the distance to the anchor. Square these nonnegative inequalities, sum, and substitute the exact RMS definition. This gives

\[\mathbb E_I\|G_I(x)\|^2\le n^2M_{\rm rms}^2\|x-x_*\|^2.\]

The integral is finite because the index space and measure are finite.



\noindent Lean declarations: \texttt{controlVariateEstimator\_second\_moment}, \texttt{controlVariateEstimator\_second\_moment\_le}.

\section{Center the error and retain the negative gradient energy}


Define the estimator error as G\_i(x)-gradient f(x). Unbiasedness proves that its uniform-index expectation is zero. Expanding its squared norm, the cross expectation is the inner product of gradient f(x) with the estimator expectation, hence the squared full-gradient norm. The exact variance identity is therefore

\[\mathbb E_I\|G_I(x)-\nabla f(x)\|^2=\mathbb E_I\|G_I(x)\|^2-\|\nabla f(x)\|^2.\]

Dropping the nonnegative final score energy gives the same n squared RMS distance bound for the centered error. Keeping it is essential for the larger-step one-step SGLD estimates.



\noindent Lean declarations: \texttt{controlVariateError}, \texttt{controlVariateError\_centered}, \texttt{controlVariateError\_second\_moment}, \texttt{controlVariateError\_second\_moment\_le}.

\section{Supply the measurable product-space error with derived L2 moments}


For a random state X with finite second moment, each component gradient evaluated at X is L2 by Lipschitz composition. A finite independent index can select these L2 functions: express the selector as the finite sum of component functions times the indicators of the measurable index events. This proves product-space L2 membership of the actual estimator and centered error. Measurability follows by checking each fixed finite-index slice. Fubini then integrates the pointwise variance bound over X, producing

\[\mathbb E\|G_I(X)-\nabla f(X)\|^2\le n^2M_{\rm rms}^2\mathbb E\|X-x_*\|^2.\]

These results supply the genuine random perturbation used by the transport comparison; no error variance is assumed.



\noindent Lean declarations: \texttt{memLp\_finite\_index\_selector}, \texttt{controlVariateEstimator\_product\_memLp}, \texttt{controlVariateError\_product\_memLp}, \texttt{controlVariateError\_product\_measurable}, \texttt{controlVariateError\_product\_second\_moment\_le}.

\section{Construct the unique zero-score minimizer}


Positive strong convexity gives a quadratic lower bound, so the sublevel set f(x) at most f(0) is bounded. Continuity and the extreme-value theorem yield a global minimizer. Fermat theorem makes its actual Frechet derivative and gradient zero. If two points had zero gradients, strong gradient monotonicity would force m times their squared distance to be at most zero. Since m is positive, the points coincide. The resulting unique minimizer provides the anchor required by the finite-sum algorithm.



\noindent Lean declarations: \texttt{exists\_unique\_minimizer}.

\end{document}
